\scanpage{4}
% archive sheet number: 3

Notons \inserted{l'isomorphisme $H^0(S,F)\simeq H^0(S'/S,F)$ \uline{déjà connu}, et} la suite exacte en basses dimensions \unsure{et|d}

\[
0\longrightarrow\check H^1(S'/S,F)\longrightarrow H^1(S,F)\longrightarrow
\check H^0(S'/S,{\unsure{\mathscr H|\mathcal H}}^1(F))\longrightarrow
\check H^2(S'/S,F)\longrightarrow H^{\unsure{0|2}}(S,F).
\]

\uline{Remarques} 1) Bien entendu, on peut énoncer Cor. 4 également pour une \uline{famille} couvrante (fpqc.)

2) La démonstration du Th. 1 montre aussi, pour $S$ affine, que

\margin{?}

\[
\check H^p(S'/S,\mathscr H^q(F))=0 \qquad\text{pour \uline{tt} }p,\ \text{si }q\ne0
\]

\deleted{Ce résultat est sans doute \unsure{ne semble pas très|vrai semble peu|la suite spectrale} \unsure{intuitif|utilisé}.}

En particulier, on trouve

\[
\check H^0(S'/S,\mathscr H^q(F))=\operatorname{Ker}\bigl(H^q(S',F)\rightrightarrows H^q(S'\times_S S',F)\bigr)=0
\]

\deleted{Il en résulte en particulier,}

\figure{\begin{tikzpicture}\draw (0,0)--(5.2,.16);\draw (1.2,-.14)--(4.45,-.07);\end{tikzpicture}}

Application à l'effectivité des Relations d'équivalence :

\uline{Proposition} \quad $X$ un préschéma, $R$ une relation d'équivalence sur \unsure{$x$|$X$}, $X_0$ sous-préschéma de $X$ défini par un Idéal nilpotent, $X_0$ stable par $R$, $R_0$ la relation d'équivalence
